\documentclass[12pt]{article}

\usepackage[margin=1in]{geometry}
\usepackage{XCharter}
\usepackage{parskip}
\usepackage{enumitem}
\usepackage{booktabs}
\usepackage{tabularx}
\usepackage{array}
\usepackage{hyperref}
\usepackage{titlesec}
\usepackage{fancyhdr}
\usepackage{needspace}

% You will submit a written update consisting of a textual description of what you have done this week and how it contributes to the progress of your project. This includes how you refined your proposal based on received feedback.

% ── Section formatting ──────────────────────────────────────────────────────
\titleformat{\section}{\fontsize{14}{16.8}\selectfont\bfseries}{}{0em}{}
\titlespacing*{\section}{0pt}{1em}{0.3em}

% ── Header / footer ─────────────────────────────────────────────────────────
\pagestyle{fancy}
\fancyhf{}
\rhead{\small Progress Report \PRNumber}
\lhead{\small \CourseName}
\rfoot{\small \thepage}
\renewcommand{\headrulewidth}{0.4pt}

% ── Fill-in macros (edit these each week) ───────────────────────────────────
\newcommand{\Section}{CS} % PUBP / CS / ECE
\newcommand{\CourseName}{CS 6727 Cybersecurity Practicum}
\newcommand{\ProjectTitle}{Multi-Signature Authentication Web Proxy}
\newcommand{\AuthorName}{Ian Barish}
\newcommand{\PRNumber}{2}
\newcommand{\ReportDate}{\today}

% --- Other Macros -----------------------------------------------------------
\newcommand{\RepoURL}{https://github.com/Ian-Costa18/Cybersecurity-Practicum/blob/progress-report-\PRNumber}
\newcommand{\DiffURL}{https://github.com/Ian-Costa18/Cybersecurity-Practicum/blob/progress-report-\PRNumber/Progress\%20Reports/diffs/diff_PR1_PR2.pdf}

% Timeline row helper (Week label, Task, Status)
\newcommand{\tlrow}[3]{#1 & #2 & #3 \\}
\newcommand{\instruction}[1]{{\small\itshape #1}}

% ────────────────────────────────────────────────────────────────────────────
\begin{document}

% ── Title block ─────────────────────────────────────────────────────────────
\noindent Section: \Section \hfill \ReportDate \\[0.2em]
\begin{center}
  {\large \textbf{\ProjectTitle}} \\[0.3em]
  \AuthorName
\end{center}

\noindent\rule{\linewidth}{0.4pt}

\vspace{0.2em}
\noindent{\small\textit{To see what has changed since the last report, \href{\DiffURL}{view the diff here}.}}
\vspace{0.4em}

% ── 1. Problem Statement ────────────────────────────────────────────────────
\section*{Problem Statement}

\instruction{Background and statement of the theoretical or practical problem, issue or opportunity motivating the research.}

Enterprise grade cryptocurrency wallets have long implemented multi-signature (``multi-sig'') authentication, which requires two or more parties to sign an outbound transaction from a secure wallet before it is sent to the receiving party. This security mechanism ensures that before a transaction is sent, it is agreed upon by the group, rather than any one individual. But what if you want to require multiple parties to approve of a critical setting change or release a package? No existing general-purpose web authentication system (think Duo or Authelia) requires multiple distinct people to cooperate before granting access.

My practicum project proposes a general-purpose multi-signature authentication web proxy that allows access to the requested resource only when multiple parties authenticate and approve the request.

Supply chain attacks have been wreaking havoc on developers in recent months. Current solutions focus on locking down individual developer accounts, but I believe a better solution is requiring multiple human verifiers to approve a project release before it goes public. My project will provide a solution that, while generalizable, will be directly applicable to this use case. This project will demonstrate this gap and then advocate for platform-native functionality to be added to the repositories.

% ── 2. Solution Statement ───────────────────────────────────────────────────
\section*{Solution Statement}

\instruction{This research will (the overall aim of the research) in order to (how the research will address the problem) with specific work product/deliverables described.}

This research will create a general-purpose multi-signature authentication web proxy that, as an example of a real-world use case, will prevent any one individual from publishing a package to PyPI and instead require multiple reviewers to approve of the package's release. I will develop an application that will be hosted on a web server. When a developer attempts to publish a package to PyPI, the web proxy will:

1. Intercept the request while saving it for publishing
2. Notify other maintainers that a package is attempting to be published
3. Gather enough approvals from the maintainers until a quorum is reached
4. Then forward the publish request on to the PyPI repository

% ── 3. Completed Tasks ──────────────────────────────────────────────────────
\section*{Completed Tasks (Last 2 Weeks)}

This Progress Report, I focused on front-loading the research, specifications, and architecture that the MVP of my project will need before committing to writing any code. I find this is ideal when working with AI to code as any knowledge it needs to create the product is all in one place. Using the \href{https://github.com/mattpocock/skills\#engineering}{/grill-me skill}, I had long conversations using Claude code to fully develop the different systems within my project, memorialized in the \href{\RepoURL/docs}{/docs} directory as specifications, Architecture Decision Records (ADRs), a Product Requirements Document (PRD), constraints the project must work in, and a threat model. These two weeks' efforts represent roughly 5,800 lines of specification created over 32 commits.

\begin{itemize}[leftmargin=1.5em, itemsep=2pt]
  % \item {[I accomplished this]}
  % \item {[and this]}
  % \item {[Discuss in greater detail than the timeline and demonstrate the
  %       expected 15 hours/week of effort]}
  \item Conducted broad cryptography research to ground the authentication scheme.
  \item Evaluated and rejected multi-signature cryptography schemes.
  \begin{enumerate}
    \item Ultimately, through research I decided to forgo complex and unnecessary multi-signature authentication cryptography solutions and instead decided on a standard credential-backed scheme for multiple reasons: (1) Multi-signature authentication schemes often require communication between users, (2) it would require every user to create and manage cryptographic keys, which is not ideal for non-technical users, and (3) it is overly complex when some of the security of the system would still rely on the security of the proxy server.
  \end{enumerate}
  \item Decided on and documented the authentication and cryptographic scheme.
  \item Defined the system architecture and request data flow.
  \item Wrote application specifications for account management, approver authentication, configuration, notification system, and the web proxy.
  \item Scoped the MVP and assessed the system's security posture through a threat model.
  \item Wrote a Product Requirements Document.
  \item Documented project decisions into ADRs.
  \item Ran a cross-document consistency audit and remediated 21 findings.
  \item Set up project process and the foundation for AI-assisted development.
\end{itemize}

% ── 4. Tasks for Next Report ────────────────────────────────────────────────
\section*{Tasks for the Next Project Report}

\instruction{What will you do over the next two weeks for your project?}

\begin{enumerate}[leftmargin=1.5em, itemsep=2pt]
  \item Decide on the tech stack for implementation of the system.
  \item Write documentation to guide agents to the documents needed to create/edit parts of the system.
  \item Break the implementation down into vertical implementation tasks.
  \item Begin writing code to implement the MVP.
\end{enumerate}

% ── 5. Questions / Issues ───────────────────────────────────────────────────
\section*{Questions I Have or Issues I'm Running Into}

\instruction{Are there unanticipated obstacles, is something taking longer than you thought it might, or is there an area where the Instructional Team might be able to provide some direction or clarification?}

I made many big decisions this Progress Report cycle that will ultimately shape how my project will come together. One example of this is the decision to use a more standard credential-backed authentication scheme rather than relying on the multi-signature cryptography schemes that I originally got the idea for this project from. However, I did not want to implement these schemes blindly just because I had sunk time and effort into researching them. I believe that a credential-based solution is ultimately simpler for the users, simpler to implement for an MVP, and just as secure as a multi-signature-based scheme. I do think that it has the potential for implementation in future versions of this project, but as a proof-of-concept to champion the idea of requiring multiple users to approve a request before it goes through, this is the solution that I want to implement. This does represent a big pivot in my thought process for designing this project, but I believe it is for the best. I would welcome the instructional team's perspective on whether taking this step forward weakens to my project's core contribution in any way, or whether the multi-party-approval concept stands on its own as a proof-of-concept.

% ── 6. Methodology ──────────────────────────────────────────────────────────
\section*{Methodology Paragraph Summary}

\instruction{Description of the methodology you are proposing to get from the starting point of your project to the final deliverable. Projects should employ some systematic process to get to your solution and a clearly identified way to evaluate it to ensure it solves the problem statement identified.}

The project will follow an iterative and research-driven development process. I will start by conducting background research on multi-signature authentication cryptography primitives, existing authentication systems, and any other relevant security design patterns. That research will then inform application specifications for the proxy and approval flow. I will then build a minimum viable product that implements the core flow of intercepting a request, collecting approvals until a quorum is reached, and forwarding the request to the target endpoint. After reviewing the MVP, I will identify any missing requirements, limitations, and security concerns, then return to do additional research and refinement of the specification before implementing the full system. I will then develop an evaluation system that ensures my project solves the intended use cases I've identified. This evaluation system will likely consist of integration tests that verify that the project works for it's intended use cases (usability), blocks attacks that attempt to bypass authentication mechanisms (security), and quantifies proxy overhead (performance), though optimizations are out of scope for this proof-of-concept.

% ── 7. Timeline ─────────────────────────────────────────────────────────────
\needspace{8\baselineskip}
\section*{Timeline}

{\small Enter one task per row. Include only project tasks (not peer feedback
or progress report submissions). Tasks must be actionable. Update the Status
column each week (Completed / In Progress / Cancelled / etc.).}

\vspace{0.5em}
\begin{tabularx}{\linewidth}{@{} l X l @{}}
  \toprule
  \textbf{Week} & \textbf{Description of Task} & \textbf{Status} \\
  \midrule
  \tlrow{W1}  {Brainstorm ideas for the project} {Complete}
  \tlrow{W1}  {Create Project Proposal} {Complete}
  \tlrow{W2}  {Conduct broad project research} {Complete}
  \tlrow{W2}  {Refine Problem/Solution Statement and Methodology} {Complete}
  \tlrow{W2}  {Define Project Timeline} {Complete}
  \tlrow{W2}  {Conduct research on multi-auth crypto primitives} {Complete}
  \tlrow{W3}  {Develop authentication scheme} {Complete}
  \tlrow{W3}  {Define high-level architecture and data flow} {Complete}
  \tlrow{W4}  {Write detailed application specifications} {Complete}
  \tlrow{W4}  {Create harness for AI assisted programming} {In Progress}
  \tlrow{W5}  {Decide on technology stack} {Not Started}
  \tlrow{W5}  {Develop minimum viable product} {Not Started}
  \tlrow{W6}  {Review MVP to identify missing requirements, limitations, and security concerns} {Not Started}
  \tlrow{W6}  {Refine application specifications based on MVP review} {Not Started}
  \tlrow{W7}  {Develop MVP evaluation framework based on use cases, security concerns, and performance} {Not Started}
  \tlrow{W8}  {Begin building fully featured application based on project specifications} {Not Started}
  \tlrow{W8}  {Create report outline} {Not Started}
  \tlrow{W9}  {Continue iterating on application development} {Not Started}
  \tlrow{W9}  {Begin writing report} {Not Started}
  \tlrow{W10} {Evaluate project based on security, usability, and performance} {Not Started}
  \tlrow{W10} {Continue writing report} {Not Started}
  \tlrow{W11} {Add finishing touches to project implementation} {Not Started}
  \tlrow{W11} {Finalize report} {Not Started}
  \bottomrule
\end{tabularx}

% ── 8. Evaluation ───────────────────────────────────────────────────────────
\section*{Evaluation}

\instruction{Include any evaluation plans and/or results by Progress Report 3. This may expand as you finalize the report.}

% ── 9. Report Outline ───────────────────────────────────────────────────────
\section*{Report Outline}

\instruction{Include an outline of your final report by Progress Report 4. This may expand as you finalize the report.}

% ── 10. References ──────────────────────────────────────────────────────────
\section*{References}

\instruction{List sources you have reviewed for your project.}

% ── 11. Appendix ────────────────────────────────────────────────────────────
\newpage
\section*{Appendix}

\instruction{If there are notes, sources, or figures you want to reference, include them here. Draft work product should begin by Progress Report 3.}

\end{document}


% Feedback

% Ian,

% Good work on PR2. For PR3, please address the professor's PR1 feedback directly, especially the need to support the motivation with evidence and to clarify the threat model and evaluation plan.

% One important writing note: the current Problem Statement still reads partly like a solution statement. Elements such as "multi-signature authentication," "multiple human verifiers," "my project will provide a solution," and "advocate for platform-native functionality" are solution or approach elements. The Problem Statement should first define the cybersecurity problem, the gap, the risk, and why the current state is insufficient. Then the Solution Statement can explain the proxy, approval flow, authentication mechanism, and implementation approach.

% Your deliverable direction is clearer now, but for PR3 please define the final deliverable more explicitly: what working MVP features will exist, what will be demonstrated, and what artifacts you will submit, such as code, documentation, threat model, evaluation tests, or a demo.

% Also, please start using clear in-text citations in IEEE or APA style. I am not deducting for citation format this time, but for PR3 and the final report, major factual or technical claims should be supported with properly formatted citations that we can follow.

% The rest of my line-level notes are in the rubric boxes and inline comments.

% You are making good progress, and the project direction is promising. Keep going, Ian. I am looking forward to seeing how your proof of concept develops in PR3, and I wish you the best of luck and success with the next stage.

% Best,
% Emad


% Comments

% Complete Template

% The report is mostly complete, but the References section is empty even though you describe research-based decisions.

% Clarity/Quality

% The Problem Statement needs to be more clearly tied to the gap being addressed. Some solution items should move to the Solution Statement.

% The report is clear overall, but key claims and design choices need stronger citation support. Without in-text citations, it is difficult to trace the evidence behind the claims.

% Please review the inline comments for line-level spelling and/or grammar notes.

% Current/Updated

% The research-to-solution link is missing:
% PR2 is meant to show a solution that is driven by background research, and that needs to be visible in the report. Each major technical and design claim should have a source or reference behind it, and you should show how the research led to the specific design choices. Please cite your references in the body at the point where they support a claim.

% Syllabus - page 7: "Project Progress Report II: Must outline solution approach and how it is informed by background research."

% Responsive

% No deduction this time. The professor's PR1 feedback was released after PR2 was due, so I am not deducting for items the student could not reasonably address before submitting PR2. However, both the PR1 feedback and this PR2 feedback should be addressed clearly in PR3.
